% =====================================================================
% 決定力場力学 論文4 — ドラフト日本語版(v0.4, 2026-08-24)
% v0.4(第191便): 論文4/5 分割(原仮定者裁定 2026-08-24「論文4/5分割: 採用」)—
%   テーマを「スケール横断の創発」から**天体再現**(マクロ: 重力・引きずり・形態)へ再構成。
%   - ミクロ章(内部スピン熱化)と隠れ質量ラダーは新設の論文5(dfm-paper5-ja.tex)へ逐語移設。
%   - 新設「太陽系実較正: ゼロフィット転写」(実測): 転写ファミリー
%     venusReal/marsMoonsReal/plutoCharonReal/uranusReal/neptuneReal+jupiterGalilean hold-out・
%     自由中心反作用則の発見(E6′-R)・引きずり超過の χ 系列 +0.002%/+0.117%/+0.246% と
%     その向き非依存(第188便)。
%   - 新設「銀河形態への段階的計画(計画)」: 8対象クラス×(宣言的解析ハローつき観測版
%     対 ダークローターの DFM 版)+事前宣言の共通ゲート。計画のみ — 結果の主張なし。
% v0.3(第122便): 論文3→論文4 へ繰り下げ(原仮定者裁定 2026-08-16 — 論文3は現実較正テーマへ
%   再割当・本稿〔スケール横断の創発〕は内容不変のまま第4論文となる)
% v0.2(第100便): 第IV節(ミクロ)・第V節(メゾ)を本文起こし(英語版 v0.2 と数値・構成同一)
% テーマ(v0.3): 「様々なスケールでの創発」(2026-08-08 採用)— v0.4 の分割で天体再現へ再構成
% 縮退経路(事前合意・v0.3): ミクロ章(第IV節)を落とし メゾ+マクロ に絞る
%   [v0.4 の分割で失効 — ミクロ章は論文5へ]
% 状態: 節構成の骨格+実測数値(tests/exp-4-84/85/89/90/91/92/93.mjs)。
%   図は全てプレースホルダ・本文は箇条書き水準で最終稿ではない。
% 英語版: dfm-paper4.tex(正本 — 数値・構成は同一)
% =====================================================================
\documentclass[a4paper,10pt,twocolumn]{ltjsarticle}

\usepackage{luatexja-fontspec}
\usepackage[haranoaji]{luatexja-preset}
\usepackage{amsmath,amssymb}
\usepackage{graphicx}
\usepackage{booktabs}
\usepackage{hyperref}

\title{決定力場力学における天体再現:\\
ゼロフィット太陽系転写・自己形成ダークローター・\\
銀河形態への段階格付け標準試験計画}
\author{(著者ブロックは論文1・2と同様 — TODO)}
\date{2026-08-24(ドラフト v0.4)}

\begin{document}
\maketitle

\section*{概要(骨格)}
% --- 主張すること ---
% 1. 一度だけ実暦に較正され(論文3)再フィットしないトイ宇宙(DFM・論文1)で、
%    マクロ側(重力・引きずり・形態)の天体再現を、共通の方法論(E0–E3 格付け+6標準試験)で進める。
% 2. 太陽系実較正(実測): ゼロフィット観測転写ファミリー(太陽金星・火星衛星・冥王星カロン・
%    天王星の環と衛星・海王星トリトン+木星ガリレオ hold-out)が kFrame=0 で観測周期を
%    10⁻⁴〜10⁻³ 水準で写す。自由中心の反作用則の穴を発見し E6′-R+ドリフト外挿自己検査で修正。
%    引きずり超過の宣言系列 χ=W/(D₀+W): +0.002%(遠)/+0.117%(χ=0.83)/+0.246%(χ≈0.92・近)—
%    公転の向きに非依存であることも実測。
% 3. メゾ: 一様な粒の雲から「重く・暗く・共回転する中心」(トイ宇宙自身の暗黒天体候補)が形成。
%    16seed・粒子数4段階・摂動・dt 4倍範囲で頑健。ソフトニング感度は正直に報告。
% 4. マクロ(応用): m=2 円盤シグネチャと四重極汚染対照。さらに連星〜棒渦巻の段階的**計画**
%    (宣言的解析ハローつき観測版 対 ダークローターの DFM 版・事前宣言ゲート)。
% --- 主張しないこと(固定リスト) ---
%    実在のダークマター/ブラックホール/銀河渦状腕の説明はしない。較正一致≠機構同定。
%    3D への定量外挿・宇宙論的存在量の主張はしない。(ミクロセクター — スピン熱力学と
%    暗い天体の光学 — は論文5。)
概要 TODO。

\section{序論}
% 問い(v0.4): 一度だけ実暦に較正されたトイ宇宙は、転写された観測から**ゼロフィット**で
% 天体系を再現できるか — 事前宣言ゲートの下で構造の階段をどこまで登れるか。
% 論文1(基礎)・論文2(箱宇宙)・論文3(現実較正 — 較正一致≠機構同定)との位置関係。
% 「1本の綺麗な動画は証拠として最下位」。正直な失敗の列挙(legacy 反作用の1公転内破壊と
% それを捕まえる自己検査・出典なし環質量のテスト質量宣言・観測に無い引きずり超過の宣言・
% 軟化長感度・大 n の整列時間尺度)。ミクロセクターは論文5へ(分割線: 論文4=重力・引きずり・
% 形態/論文5=スピン・熱・光学)。
TODO。

\section{構造形成に関わるモデル追加(最小再掲)}
% W・重力 a=G∇W・フレーム引きずり(A8/E6′)・内部温度 T_int(A9″)・
% C·m 重み伝導(E10′)・融合+コア継承・コアv2(J_core)・減光 lS_eff(A11/E8R)・
% E6′-R 対称反作用(支配天体を固定しない閉鎖系)。導出は論文1へ委譲。
TODO。

\section{創発の標準(方法論の中核)}
% 表1: E0–E3 の定義。6標準試験の定義と運用規則(秩序変数の事前宣言・
% 負の結果も同一書式で記録)。
TODO。

\section{太陽系実較正: ゼロフィット転写}
% [第191便: 実測節。以下の数値はすべてコミット済みの値 — 内蔵プリセット venusReal/
%  marsMoonsReal/plutoCharonReal/uranusReal/neptuneReal(beta/index.html)・出典表
%  paper/data/solar-observations.csv(全行アーカイブURLつき)・QA behavior.* / ai.obs-*。
%  本文は第一次パス(英語版 dfm-paper4.tex 第IV節が正本 — 数値・構成同一)]

\subsection{転写の契約}

銀河を試みる前に、較正済みセクターは「フィットしたことのない系」との接触を生き延びねば
ならない。アプリは決定的な転写経路を持つ: 観測レコードの表(1行=1観測量・出典と実際に
開いた URL つき)が、宣言された決定的関数\emph{だけ}によって走る2次元系へ変換される。
サンプル別の実単位はスケール指数の1径数族 $L-T=4$・$M+2T-3L=11$(実定数 $G$・$c_0$・
$\kappa=G/c_0^2$ を保つ唯一の組)として宣言し、qLock 減衰 $q$ は厳密一致式
$q_{\mathrm{exact}}$ を宣言参照軌道で1回だけ評価した直値、共有補正 $D_0=0.006$ は論文3の
月較正からの流用(\emph{系ごとの再フィットはゼロ})、初速は較正係数ちょうど $1.000$ の実ケプラー
速度である。各構築は自己診断(400步の試走で推定周期を転写周期と $\pm1\%$ で照合)を通らねば
採用されない。理想化 — 2次元赤道面・離心率が出典に無い場合の円軌道化・値域下限を割る質量の
テスト粒子降格 — はサンプルごとに宣言され、黙って行われることはない。

\subsection{ファミリーと、その実測}

内蔵の機械検査つきサンプルとして6系を持つ: 木星ガリレオ系(論文3の hold-out)・太陽金星・
火星とフォボス/ダイモス(宣言つきテスト粒子転写)・冥王星カロン(重心が主星の外にある二体)・
天王星の環と主要5衛星・海王星トリトン(逆行衛星の二体)。実測の要点は3つある。

\emph{構造的発見。}自由中心の二体転写を、固定された支配天体を反作用リザーバと前提する
legacy 反作用則のまま走らせると、$\chi=W/(D_0+W)\rightarrow 1$ の系で軌道は1公転以内に
壊れる(太陽金星の実測: 半径振れ幅 45\%・周回不能)。修正はフィットではなく構造である:
二体重心系を対称化反作用則 E6$'$-R で構築し、自己診断に\emph{ドリフト外挿検査}(二体限定 —
試走窓の周期ドリフトを1公転へ外挿して差し戻す)を加えた。legacy 構築は外挿 $1.47\%$ で
拒否される(この否定対照自体が QA の一部である)。則の修正後、金星年は観測 224.701 日に
対し 224.706 日($+0.0020\%$)に載る。

\emph{宣言される引きずり超過と、その $\chi$ 系列。}$k_{\mathrm{frame}}=0$ では転写は
観測周期を $10^{-4}$〜$10^{-3}$ 水準で写す(海王星トリトン: 観測 5.876854 日に対し
5.87741 日・$+0.009\%$)。雛形規約 $k_{\mathrm{frame}}=1$ では運動引きずり項 E6$'$ が
観測に無い量だけ周期を押し上げる — これを吸収せず宣言する: 太陽金星 $+0.002\%$
($\chi\approx1.0$ だが遠距離)・海王星トリトン $+0.117\%$($\chi=0.83$)・冥王星カロン
$+0.246\%$($\chi\approx0.92$・近接)— $\chi$ と軌道の近さの両方が効く系列である。
この系列が、銀河計画(計画節)が正直に向き合うべき「引きずり対観測」の張力の実測基盤になる。

\emph{向き非依存。}海王星トリトンは実在どおり\emph{逆行}で構築される: $\pi/2$ を超える軌道
傾斜(トリトン: 出典表から転写した $157.345^{\circ}=2.746188$ rad)は公転の向きとして
決定的に転写され、接線速度の符号だけを反転する(面外の傾きは無視して宣言)。逆行版の周期・
半径振れ幅は、両方の $k_{\mathrm{frame}}$ 設定で順行版と測定精度内で同一である: この系の
引きずり超過は公転の向きに依存しない。機構帰属はせず、記録する。

正直な宣言が2つ残る: 天王星の環質量の出典はモデル依存の推定のみ(掩蔽自己重力モデル —
ε 環で $3.27$〜$6.99\times10^{15}$ kg)で、粒あたり $\sim10^{-10}$ 単位と値域下限未満のため、
環粒子は宣言つきの下限値で置かれる(半径構造は転写され、実測窓で $k_{\mathrm{frame}}=1$
は 0.0164 単位・$k_{\mathrm{frame}}=0$ 対照は 0.968 単位の保持 — 記録のみ・機構帰属せず)。
転写経路は構成上2次元で、軌道傾斜は向きとしてだけ入る。

\subsection{本節の主張}

主張は意図的に狭い: 一度だけ較正された弱場セクター+決定的転写契約は、系ごとのフィット
ゼロで太陽系の周期を $10^{-4}$〜$10^{-3}$ 水準で再現し、残る不一致は宣言され・定量され・
名前のある項へ遡る。較正一致は機構同定ではない(論文3)— 実在の太陽系がモデルの引きずり項を
持つとは、ここでは何も主張しない。

\section{メゾスケール: 自己形成ダークローター(中心結果)}
% 実測(exp-4-84/85/89/92): 180粒→融合155回→中心天体 26.7%(対照 0.56% 厳密)/
% G 用量反応 単調/ J_core 継承恒等式(48個ぶん=飲み込んだ質量と1粒単位で一致)/
% 16seed: mFrac 18.3〜53.3%・減光コントラスト 2.41〜5.30 vs 対照 0.40〜1.26(重なりゼロ)/
% dt 0.016/0.008/0.004 収束(秩序変数が seed 幅の内側・relL 1e-6 台)/
% ε×0.5〜×2: 形成は残るが mFrac は ε×2 で下振れ(正直な感度報告 — 定量は ε 固定で引用)/
% n=90/180/360/720 の4段階: n=720 で形成完全(対照比166〜208倍)・mFrac 飽和・
% 整列時間尺度が n とともに延びる(9000步 0.752 → 13500步 1.000 — 時間窓の注記)。
% 立たなかった主張(暗さ単独は判別しない・kFrame 対照の系統差なし・dFrac 非単調)も収載。
% [第100便: 本文起こし — 数値は第97便の c₀=30 単位系での全面再実測を正とする。
%  exp-4-89/92(dt/ε/n=720)のみ再パラメータ化前の旧単位記録 — 本文で明示。
%  英語版 dfm-paper4.tex(旧 dfm-paper3.tex v0.2)の対応節と数値・構成同一]

\subsection{系}

メゾスケールの系は、同一の粒 180 個(質量1・スピン0・全て自由)の
一様円盤に、遅い全体回転($\omega_0$)と、全粒に同一の微小コア種
(質量の5\%・$J_{\mathrm{core}}=0.0010$ ずつ)を与えたものである。
それ以外は何も置かない: 中心天体も、暗いコアも、ハローも、固定粒子も、
境界も、外場もない。働く規則は自己重力(E4)・接触(E9)・融合
(2粒の中心間距離が $d<0.7\,(R_i+R_j)$ で1粒になる)・コアの保存則
継承($M_c' = M_{c,i}+M_{c,j}$, $R_c' = \sqrt{R_{c,i}^2+R_{c,j}^2}$,
$J_c' = J_{c,i}+J_{c,j}$)である。重心と重心速度は build 時に一度
だけゼロ化し、以後は閉鎖系(E6$'$-R)— 保存量モニタは運動量・
角運動量を全リザーバ簿ゼロで閉じる(基準 run で
$|\Delta L|/L_{\mathrm{scale}} = 3.5\times10^{-6}$・16 seed 全体で
$\leq 1.4\times10^{-5}$)。

以下の数値はアプリの一律光速規約($c_0=30$)の単位系によるもので、
本サンプルはこの規約で全面再実測済みである。基準 run は
seed 20260806・$t=288$(18000 步)。

\subsection{形成と、その原因}

基準 run では 180 粒が 146 回の融合を経て 34 体になり、最重の天体は
全質量の 27.2\%($m=49$)を持つ中心へ育つ(束縛率 0.94)。
ノックアウト対照(融合オフ)では最大天体は元の1粒のまま — 質量比は
全 seed で厳密に $0.56\% = 1/180$ — 束縛率は 0.15: 重い中心は
融合則の産物であり、初期条件の産物ではない。自己重力の因果は用量反応で
分離する: $G$ だけを $0.5/1/1.5/2/3$ と変えると質量比は
$1.1/7.8/18.9/27.2/57.8\%$($7/62/125/155/173$ 融合)と単調に動く。

コア角運動量は継承経路だけで蓄積する: 中心の
$J_{\mathrm{core}}=0.0490$ はちょうど種 49 個ぶんで、飲み込んだ質量
49 と1粒単位で一致する(継承則の和が文字どおり数えている)。
系全体の $\sum J_{\mathrm{core}}$ は最後まで
$0.1800 = 180\times0.0010$ に留まり、コア種を除いた対照では
$J_{\mathrm{core}}\equiv 0$ である。

暗さも育つ。実効減光(A11/E8R — 天体が自転で自光をどれだけ掻き出すか)
は全粒 $0.000$ から始まり、$t=288$ で中心天体は $0.818$ まで暗くなる
一方、周囲の平均は $0.192$ — 暗い中心と明るい周囲、コントラスト 4.3
が生成される。減光は井戸の深さ $\psi=W\kappa$($\kappa$ は論文1の
時空係数 — 一般相対論のアインシュタイン重力定数 $8\pi G/c^4$ とは別の量で、
スピン熱伝導 $\kappa_s$ とも異なる)だけで決まり力学には
一切入らない: $\kappa$ を $1/12,\,1/14,\,1/16,\,1/20,\,1/24$ と掃引しても粒子数・融合数・
質量比・スピン・半径は最後の1ビットまで同一で、減光だけが
$1.000/0.818/0.574/0.349/0.251$ と動く。粒はスピン0から出発して
全会一致の共回転に至る(整列 $0.000\rightarrow1.000$・クランプ数
$166\rightarrow2$・最大クランプ質量比 0.994・回転支持
$V/\sigma=11.2$)。

\subsection{6つの標準試験}

\emph{多seed(試験3)。}16 seed で中心質量比は 16.7〜47.8\%
(中央値 26.7\%)・全 seed で 16.7\% 以上。融合なし対照は全 seed で
0.56\% に張り付く — 最悪 seed でも対照の 30 倍である。同じ 16 seed で:
最大クランプ質量比 0.978〜1.000・整列は全 seed で 1.000・$V/\sigma$
4.1〜22.0・中心減光 0.513〜1.000 対 周囲 0.156〜0.220・
$J_{\mathrm{core}}$ 0.0300〜0.0860(種 30〜86 個ぶん)・束縛率
0.49〜0.94・残り 26〜42 体(138〜154 融合)・NaN なし・クランプなし。
光学量の判別子は\emph{コントラスト}(中心/周囲)である: 本則
2.59〜5.41(中央値 4.11)対 対照 0.36〜1.26(中央値 1.16)—
両レンジは一度も重ならない。

\emph{ノックアウト・用量(試験1–2)}は上の融合ノックアウト・コアなし
対照・$G$ 用量である。融合しきい値 $d_{\mathrm{frac}}$ の第二の用量
反応($0.7/0.6/0.5/0.35 \rightarrow 27.2/37.2/53.9/53.9\%$)は
\emph{この実現では}単調→飽和だが、前のパラメータ化では同じ掃引が
非単調($26.7/47.8/43.3/23.3\%$)と実測された。この用量曲線の形
そのものが丸めの実現に敏感であり、きれいな方の曲線を選ぶのではなく
その事実を記録する。

\emph{粒子数スケーリング(試験4)。}$n=90/180/360$(代表3 seed)で
質量比の中央値は 17.8/22.8/37.5\%・最大クランプ比 0.933/0.994/1.000・
整列は全条件 1.000。正直に述べるべき絡み合いがある: 本サンプルは
円盤半径と粒質量を固定したまま $n$ を変えるので、$n$ は総質量
(=自己重力の強さ)でもある — $n=90$ は $G$ 半減に相当し、対照10倍
の基準をそこで満たすのは 3 seed 中 2 つ(比 2.0/16.0/20.0)である。

\emph{摂動回復(試験5)。}形成後の $t=240$ に、アプリ正式の運動量
保存injectorで全粒へガウス速度ノイズを注入する。速度RMSの30\%では
どの秩序変数も動かない。振幅を掃引すると最大クランプが壊れ始めるのは
速度RMSの4倍から($0.994\rightarrow0.861$。2倍までは 0.994 のまま)。
乱されて戻るのは $V/\sigma$(39〜61\% 低下 → 20〜100 步で復帰)と
束縛率($\leq$15\% 低下 → 20〜50 步で復帰。1 seed は測定窓内で回復
基準を外す — 記録)。速度ノイズは定義上整列に触れないので、別の摂動
として約半数の粒のスピンを反転する: 整列は 1.000 から
$0.531/-0.212/0.084$ へ落ち、20 步($\Delta t=0.32$)で復帰する。
共回転状態は一回きりの配置ではなく、接触摩擦(E9)と引きずり
(A8/E6$'$)が毎步立て直すアトラクタである。

\emph{時間窓(試験6)。}クランプは 3000 步($t=48$)までに1つへ
併合し、整列は 4000〜5000 步で 0.99 を超え、質量比は 4000〜6000 步で
最終値の90\%に達して飽和する — 宣言窓の 18000 步と 36000 步の差は
高々 0.6 ポイント。窓は立ち上がりを過ぎた平坦域の内側にある。

\subsection{数値頑健性と世代注記}

さらに3つの掃引は物理ではなく積分を調べる。これらは本サンプルの
\emph{再パラメータ化前}(\texttt{tests/exp-4-89.mjs}・
\texttt{exp-4-92.mjs}、2026-08-08)の記録である — その後、サンプルは
一律光速規約 $c_0=30$ への厳密相似変換(速度・率 $\times k$・時間
$\div k$)で再パラメータ化された。連続力学は構成上同一・離散積分は
僅かに異なり、上の多seed分布は再パラメータ化後に再実測済みである。
以下は旧単位系のまま、旧単位系の16seed下限 18.3\% に対して引用する:

(i)\emph{時間刻み}: $dt \in \{0.016, 0.008, 0.004\}$(物理時間固定)
で全秩序変数が16seed幅の内側・
$|\Delta L|/L_{\mathrm{scale}}\sim10^{-6}$。融合数は緩やかにドリフト
($155\rightarrow147\rightarrow138$)。(ii)\emph{軟化長}:
$\times0.5/\times1/\times2$ — 形成はどこでも残る(最大クランプ
0.978〜1.000・整列 1.000)が、質量比は $\times2$ で\emph{確かに}
軟化長に敏感(融合 104〜122 に減・最小質量比 11.1\% は16seed下限を
割る): 定量的な質量比は軟化長固定で引用しなければならず、本稿も
そうしている。(iii)\emph{$n=720$}: 形成は完全(最大クランプ 1.000・
対照比 166〜208 倍)・質量比は飽和(中央値 31.5\% 対 $n=360$ の
31.7\%)だが、整列時間尺度は $n$ とともに延びる — ある seed は
9000 步で 0.752、13500 步で 1.000。$t$ 窓は $n=180$ で調整された
ものであり、$n=720$ では注意書きとなる。そのとおりに記録する。

\subsection{メゾ章の主張}

一様な同一粒の雲から、重く・暗く・共回転し・角運動量を持つ中心が
閉鎖系の中で形成される。帰結の全要素はノックアウトか用量反応で宣言
された規則に遡り、16 seed と3つの粒子数で再現し、摂動から回復し、
2つの正直な感度(軟化長・大 $n$ での整列時間尺度)を明記して述べ
られる。第III節の格付けでこれは E3 — 頑健な創発である。ただし創発
するのは\emph{このトイ宇宙自身のダークローター}であって、実在の空の
何ものでもない(第IX節)。

\section{マクロスケール: 円盤シグネチャ(応用)}
% (v0.4: 隠れ質量ラダーは分割線〔光学・簿記=ミクロ側〕に従い論文5へ移設)
% m=2 シグネチャの spin 用量反応と四重極汚染対照。
% 語法規則: 「銀河を説明」ではなく「トイモデル内でどんなシグネチャを刻むか」。
TODO。

\section{銀河形態への段階的計画(計画)}
% [第191便: 計画節 — 結果の主張は無い。全ゲートを測る前に宣言する]

本節は計画であり、測る前にゲートを宣言し、結果を一切主張しない。対象は構造の階段順に
8クラス: 普通星連星・連星パルサー・ブラックホール連星・球状星団(判別実験として矮小楕円体
銀河を挿入)・楕円・レンズ状・渦巻・棒渦巻。各クラスは転写観測から\emph{2通り}構築する:
そのクラスの標準解析が要求する場合に宣言的な解析ダークハロー項を含む\emph{観測再現版}
(ハローのノックアウトが対照。連星はハローが期待されないスケールなので陰性対照を担う)と、
ハローをトイ宇宙自身のダークローター(メゾ節)で置き換え同じ標準で格付けする \emph{DFM 版}
である。

全段階は事前に固定した1つのゲート一式を共有する: NaN・クランプゼロ/保存則帳簿/時間刻みと
軟化長の半減収束/多seed再現/摂動回復/時間窓の宣言/ノックアウト対照(ハロー・ローター・
$k_{\mathrm{frame}}$)/駆動主張には単調用量反応/凍結後にのみ判定する保留データ/全公表値への
4値主張ラベル。\emph{しないこと}も今宣言する: ブラックホール連星は合体前で止め波形マッチを
しない。連星パルサーの軌道減衰は観測版の\emph{外部簿記項}としてだけ持ち、エンジン公理には
決して入れない。対のローターで駆動した棒・腕の応答は強制 $m{=}2$ 対照とラベルし、自発形成と
呼ばない。共有補正 $D_0$ を銀河ごとに再フィットしない。

段階1の前に3つの前提整備が要る: 対等質量連星の一般化観測スキーマ(E6$'$-R の自由中心 —
その必要性は太陽系節が実測で立てた)・観測版のための解析ハローキット・サンプル別パラメータ
監査を拡張する誤差/保留データ宣言。太陽系節の引きずり超過系列は、全段階が自分の観測に対して
報告すべき常設の張力として持ち越される。

\section{数値頑健性のまとめ}
% 現象×(dt・ε・n・seed・時間窓)の横断表。例外(ε 感度・整列時間尺度・伝導時間尺度)を
% 強調表示 — 隠さない。
TODO。

\section{限界と負の主張}
% 固定リスト(v0.4 改訂): 2D/order-unity引きずり/Lorentz不変なし/因果伝播なし/減光は
% モデル固有/融合は粗視化/ダークローター≠ダークマター/DFM BH≠GR BH/存在量主張なし/
% 転写は2D赤道面(π/2 超の傾斜は逆行の向きとしてだけ転写 — 宣言)/較正一致≠機構同定
% (引きずり超過系列は宣言であって実在太陽系の引きずりの証拠ではない)。
% (熱アナロジーの限界項はミクロセクターとともに論文5へ。)
TODO。

\section{再現性}
% 現象→プリセット→ハーネス→seed集合→時間窓→コミット→結果JSON の台帳表。
% ハーネス: exp-4-84(ノックアウト/用量)・4-85(多seed/N/摂動/時間窓)・4-89(dt/ε)・
% 4-92(n=720)。太陽系節: 転写経路 QA(ai.obs-*)・サンプル別 QA(behavior.venusReal/
% marsMoonsReal/plutoCharonReal/uranusReal/neptuneReal)・出典表
% paper/data/solar-observations.csv(全行アーカイブURL)。
% (exp-4-90/91/93〔ミクロ〕は論文5から引用。)
TODO。

\begin{thebibliography}{9}
\bibitem{paper1} 論文1(TODO)
\bibitem{paper2} 論文2(TODO)
\bibitem{paper3} 論文3(現実較正 — TODO)
\bibitem{paper5} 論文5(ミクロスケールの相変移 — TODO)
\end{thebibliography}

\end{document}
